\section{Primitive Heisenberg carriers without tensor indecomposability}
\label{sec:heisenberg-primitive-tensor-carrier}

\paragraph{Author freeze and provenance.}
D-TD38, reserved question Q2et, is a separate extension of D-TD36.
Root proposed terminal tensor factorization, the two dimension ranges
and the larger natural-degree threshold.  The argument below supplies
the actual coordinate-quotient, geometric-rank and exceptional-core
interfaces in full.  It uses the separately reviewed commuting-field
bound D-TD34 and the accepted arbitrary-partner source obstruction
HD-FT27.  The complete author proof, source-interface, checker and
first-output readings pass.  Independent and root whole-proof reviews also pass.
No normal subgroup of the original group is inferred from an
abstract alternating quotient.

\paragraph{Actual setup.}
Use the actual setup (HCF.1)--(HCF.3) of D-TD34: an actual connected
finite $J$-cover of the sphere, a finite coefficient representation
$V=W\otimes R$ with its entire external partner, and
\[
 \begin{gathered}
 e=\dim V,\qquad \sum_i\operatorname{rank}(1-V(g_i))=2e+6,\\
 P=\operatorname{im}(J\to\operatorname{PGL}(W)),\qquad
 Q\triangleleft P,\qquad Q\simeq\mathbf F_p^{2m},\\
 H=P/Q\hookrightarrow\operatorname{Sp}(Q),\qquad
 \psi:H\twoheadrightarrow T\in\{A_n,S_n\}.
 \end{gathered}
 \tag{HPT.1}
\]
The projective label action on $W$ is nondegenerate and irreducible,
and the actual defining $\mathbf F_pH$-module $Q$ is irreducible.
Its intrinsic commuting field and absolute model are
\[
 F=\operatorname{End}_{\mathbf F_pH}(Q)=\mathbf F_q,\qquad
 q=p^f,\qquad a=\dim_F Q,\qquad af=2m,\qquad
 \dim_{\mathbf C}W=q^{a/2}.
 \tag{HPT.2}
\]
The finite-algebra proof in D-TD34 makes this $F$-model absolutely
irreducible.  Primitivity below is an additional hypothesis.

\begin{theorem}[D-TD38]
Suppose that $n\ge2^{30}$ and that the $a$-dimensional $F$-model
of $H$ is primitive after extension to an algebraic closure of $F$.
Then the actual source (HPT.1) is impossible.
No tensor-indecomposability hypothesis is imposed on that model.
No absolute irreducibility of the original $\mathbf F_p$-module,
bound on $\dim R$, or support floor on $R$ is assumed.
\end{theorem}

\begin{proof}
\paragraph{Step 1: the field bound and the whole absolute dimension.}
The hypotheses of D-TD34 are satisfied, so $q<28900<2^{15}$.
In particular the actual faithful matrix group $H$ is defined over
this bounded field.  If $a<8192$, then
\[
 |H|<q^{a^2}<2^{15\cdot2^{26}}
       <2^{29\cdot2^{29}-1}
       <\frac{(2^{30})!}{2}\le\frac{n!}{2}\le|T|.
 \tag{HPT.3}
\]
The penultimate strict comparison follows by taking the last
$2^{29}$ factors of $(2^{30})!$, each strictly larger than $2^{29}$.
The exponent comparison is $15\cdot2^{26}<29\cdot2^{29}-1$.
Surjectivity in (HPT.1) contradicts (HPT.3).  Therefore
\[
                            a\ge8192.
 \tag{HPT.4}
\]
Only the whole original matrix image is bounded by $q^{a^2}$.
No field bound on a chosen terminal factor will be presumed.

\paragraph{Step 2: finite terminal factor images over the algebraic closure.}
Put $\overline F=\overline{\mathbf F}_p$ and
$M=Q\otimes_F\overline F$.  If $M$ has a nontrivial tensor
decomposition whose two full factor matrix algebras are separately
normalized by $H$, split it.  Repeat on the resulting coordinate
projective images.  Dimensions strictly decrease at each split,
so this produces a finite factorization
\[
 M=M_1\otimes\cdots\otimes M_s,\qquad
 r_i=\dim M_i\ge2,\qquad a=\prod_i r_i,\qquad
 P_i=\operatorname{im}(H\to\operatorname{PGL}(M_i)).
 \tag{HPT.5}
\]
Here $s=1$ is allowed.  Every $P_i$ is an actual finite image,
and $H$ normalizes each displayed factor algebra.  Later splits
are made in an actual coordinate image and pulled back to $H$;
there is no residual multiplicity action outside (HPT.5).

Each factor is irreducible: a proper invariant subspace of one
factor, tensored with all the others, would be $H$-invariant in
$M$.  Each factor is primitive: an imprimitivity direct sum in
one factor would similarly give an $H$-permuted direct sum in
$M$.  By termination every factor is tensor-indecomposable in
the precise Guralnick--Saxl convention.  Only factorwise preserved
tensor products have been split; permutations of equal tensor
factors are not silently excluded by that convention.

For each $P_i$, its full inverse image in
$\operatorname{SL}(M_i)$ is finite.  Indeed each projective
matrix class has at most $r_i$ determinant-one representatives
over the algebraic closure, and the scalar kernel is finite,
also when $p$ divides $r_i$.  Denote this finite linear group by
$G_i$.  It acts irreducibly, primitively and tensor-indecomposably
on $M_i$.  The entries of these finite matrices may lie in a
larger finite extension of $F$; nothing requires (HPT.5) or $G_i$
to descend to $F$.

Absolute irreducibility gives
$Z(H)=H\cap F^*$.  This center maps trivially to the centerless
$T$.  The combined coordinate map has kernel exactly $Z(H)$:
if every factor operator is scalar, its tensor is scalar, and
the converse follows by centralizing the full factor matrix
algebras.  Thus there is an actual embedding and surjection
\[
 \overline H:=H/Z(H)\hookrightarrow\prod_{i=1}^s P_i,
                 \qquad \overline\psi:\overline H\twoheadrightarrow T,
 \tag{HPT.6}
\]
with every coordinate projection onto $P_i$.

\paragraph{Step 3: the actual alternating quotient descends to one coordinate.}
Write $K_i$ for the kernel of the $i$th coordinate in (HPT.6).
These are normal subgroups with $\bigcap_iK_i=1$.
Every nontrivial normal subgroup of $T=A_n$ or $S_n$ contains
$A_n$, in our range of $n$.  Suppose for contradiction that
$\overline\psi(K_i)\ne1$ for every $i$.
For $I_j=\bigcap_{i\le j}K_i$, induction gives
$\overline\psi(I_j)\supseteq A_n$: the initial case holds,
and
\[
 [I_{j-1},K_j]\le I_j,\qquad
 \overline\psi([I_{j-1},K_j])
   =[\overline\psi(I_{j-1}),\overline\psi(K_j)]\supseteq A_n.
 \tag{HPT.7}
\]
The last inclusion uses $[A_n,A_n]=A_n$.
This contradicts $I_s=1$.  For some $i$ the quotient therefore
kills $K_i$, and factors as an actual surjection
\[
                     \psi_i:P_i\twoheadrightarrow T,
              \qquad \psi=\psi_i\circ\pi_i.
 \tag{HPT.8}
\]
The same conclusion for $s=1$ is immediate.  This is the elementary
coordinate-kernel argument of D-TD17, supplied here in full.
It is stronger than merely locating a composition factor.
Fix this coordinate and put $r=r_i$, $G=G_i$.
Composition with $G\to P_i$ supplies an actual surjection
$G\to T$, also denoted by $\psi_i$.  Its scalar center maps
to one; the whole kernel need not be central.

\paragraph{Step 4: geometric rank transfers to the whole original module.}
For a matrix $B$ over $\overline F$ put
$\nu_{\rm alg}(B)=\dim B-\max_\lambda\dim\ker(B-\lambda I)$.
For arbitrary invertible matrices $A,B$ in every characteristic,
the complete Jordan-intertwiner argument of D-TD27 proves
$\operatorname{rank}(A\otimes B-\lambda I)
\ge(\dim A)\nu_{\rm alg}(B)$ for every scalar $\lambda$.
Indeed matching Jordan blocks of sizes $u,v$ contribute
$\min(u,v)$ to the intertwiner space; summing bounds the kernel
by $\dim A$ times the largest geometric eigenspace of $B$.

For an actual $h\in H$, choose a representative $B_h\in G$
of its $i$th projective coordinate.  Absorb any reciprocal scalar
in the tensor of all other coordinates.  Rank is unchanged on
extending the original $F$-matrix to $\overline F$ and is
multiplied by $f$ on restricting scalars.  Consequently
\[
 \operatorname{rank}_{\mathbf F_p}(1-h)
   =f\operatorname{rank}_{\overline F}(1-h\mid M)
   \ge f\frac ar\nu_{\rm alg}(B_h).
 \tag{HPT.9}
\]
The bound works even for nonsemisimple $h$, and for every possible
scalar representative of its coordinate.  If $\psi(h)\ne1$,
then $B_h$ is nonscalar by (HPT.8), so its geometric projective
rank is at least one.

If $r\le575$, (HPT.4) and (HPT.9) give original prime-field
rank at least $8192/575>14$.  This already gives the uniform
rank bound at least thirteen which will be used below.
Suppose instead that $r\ge576$.  Apply the exact published
Guralnick--Saxl Theorem7.1 and Remark3 to the terminal finite
linear group $G$ over $\overline F$.\cite{BI:GS}
In its generic branch every nonscalar coordinate has geometric
projective rank strictly greater than $\sqrt r/2$, hence integer
rank at least thirteen.  Again (HPT.9) gives original rank at
least thirteen for every quotient-nonidentity $h$.

In either of these two branches the exact D-TD05 Heisenberg
trace formula gives normalized trace at most
\[
                    c=2^{-13/2}<\frac1{90}.
 \tag{HPT.10}
\]
For completeness retain the full original periods in $J$.
If an original element has order $L$ and natural quotient order
$o$, exactly $L/o$ of its powers have trivial natural image.
Their normalized whole-coefficient traces are at most one;
the others are at most $c$, with the entire external partner
retained by reciprocal finite unitary tensor normalization.
Cyclic Fourier inversion therefore gives moving rank at least
$e(1-c)(1-1/o)$.  The actual connected $T$-quotient curve has
genus at least two and orbifold sum at least $85/42$, by the
complete actual-cover argument of D-TD02.  Summing yields
\[
                 e(1-85c)\le252,
                 \qquad e<4536,
                 \qquad e\ge q^{a/2}\ge2^{4096}>4536,
 \tag{HPT.11}
\]
a contradiction.  No cover of the artificial lift $G$ has been
introduced: the cover here is the actual quotient of (HPT.1).

\paragraph{Step 5: the terminal exceptional core.}
It remains to consider a terminal coordinate of degree $r\ge576$
in a Guralnick--Saxl exceptional branch.  Since $r>10$, only
the full natural classical and the deleted alternating/symmetric
cases remain.  The supplied Section7 proof identifies
$S=F^*(G/Z(G))$ as simple in these cases.  Let $D$ be its central
preimage and $E$ the stable derived subgroup of $D$.
It is an actual normal perfect subgroup of $G$ mapping onto $S$
with central kernel, hence is quasisimple.

The nonzero $E$-socle of $M_i$ is $G$-invariant, so equals the
irreducible module.  Primitivity makes its restriction homogeneous.
The irreducible $E$-factor has dimension greater than one, and
its evaluation tensor product has both full factor algebras
normalized by $G$.  Terminal tensor indecomposability forces
multiplicity one.  Thus $E$ is absolutely irreducible on $M_i$,
and its centralizer in $G$ is the scalar center.

An automorphism of perfect $E$ which induces the identity on
$E/Z(E)$ is the identity: its difference from the identity is
a homomorphism into the center.  Therefore
$G/(EZ(G))\hookrightarrow\operatorname{Out}(S)$.
The exact published Schreier statement makes $G/E$ solvable.
\cite[Theorem2.3, p.~1314]{HD:BabaiLuksSeress1997}
Under the actual surjection $G\to T$, the normal perfect image
of $E$ is nontrivial, because $T$ is nonsolvable, and is therefore
$A_n$.  It follows that $S\simeq A_n$.
The full classical/alternating repetition list in Li--Xia,
Section2.2.1, numbered author p.13, permits only $A_5,A_6,A_8$
in the named natural classical families.\cite{HES:LiXia2022}
It excludes this full classical core for the present $n$.
No assertion about alternating subgroups of classical groups
is made or needed.

\paragraph{Step 6: a deleted terminal coordinate, with all scalar twists.}
In the remaining case the actual perfect $E$ is the ordinary
$A_c$ in its fully deleted module.  The derived group of a
scalar extension of an ordinary perfect group is that group
itself, so passing to the finite special-linear lift introduces
no different perfect central cover.  Its simple image $A_n$
gives $c=n$, and the map on $E$ is injective.  Relabeling by
an automorphism of $A_n$, which is natural in this range,
identifies the coordinate module over $\overline F$ with
$D_n(\overline F)$ and its actual quotient with the prescribed
natural letters.  Its dimension is $r=n-1$ or $n-2$.

For every $g\in G$, the represented operator and
$D_{\psi_i(g)}$ implement the same conjugation on this
irreducible ordinary core.  Schur's lemma makes their ratio
a scalar in $\overline F^*$.  Accordingly each $B_h$ above
has the form $\lambda_hD_{\psi(h)}$.  This identifies every
actual coordinate operator, not just the elements of its core.
No descent of this tensor coordinate to $F$, splitting of a
central extension, or identification of an invariant form is needed.

The elementary cycle equations of D-TD35 apply in every
characteristic and for every scalar twist.  They give
\[
 \nu_{\rm alg}(B_h)\ge\operatorname{ind}_n(\psi(h))-2,
 \qquad \psi(h)\ne1\ \Longrightarrow\quad
                       \nu_{\rm alg}(B_h)\ge1.
 \tag{HPT.12}
\]
For the first inequality, each cycle of a scalar times a
permutation has kernel dimension at most one; passing to the
augmentation space and its possible constant quotient loses
at most two from rank.  This holds for every scalar target,
so it bounds geometric projective rank.  For the second,
projective faithfulness follows by applying a putative scalar
deleted operator to each $e_i-e_j$: a constant error has
support at most four, and hence vanishes for $n\ge5$.
Thus the natural permutation fixes every unordered pair and is one.

Use (HPT.9), with $f a/r\ge1$.  If a nonidentity natural
operator moves $k\le n/16$ letters, its Heisenberg trace ratio
is at most $1/\sqrt2<(15/16)^3\le((n-k)/n)^3$.
If $k>n/16$, then its ordinary prime-field rank is greater
than $n/32-2$, and its trace ratio is less than
$2^{1-n/64}<n^{-3}$.  The last dyadic bound is the elementary
D-TD22 estimate, valid for $n\ge4096$.  Also
\[
 e\ge q^{a/2}\ge2^{(n-2)/2}>n^3,
 \qquad
 \frac{|\operatorname{tr}_V(g)|}{e}
       \le\left(\frac{\max\{\operatorname{fix}_n(\psi(g)),1\}}n\right)^3,
 \qquad
 \psi(g)\ne1\ \Longrightarrow\quad
       \frac{|\operatorname{tr}_V(g)|}{e}<\frac78.
 \tag{HPT.13}
\]
Here $\psi(g)$ denotes the actual composition from $J$ in
(HPT.1).  Reciprocal unitary factor lifts retain all original
powers and the entire external partner exactly as in D-TD35.
At natural image one the unitary bound one suffices.

Finally pass to the actual faithful coefficient image if necessary.
An element of $\ker V$ has scalar $W$-factor, since a tensor
operator equal to the identity has scalar factors.  It therefore
kills $P,H,T$.  The connected cover $C/\ker V$ carries the
descended actual $T$-quotient and the same moving-rank sum;
identity coefficient inertia images contribute zero.
All bounds (HPT.13) descend.  Accepted HD-FT27 now supplies
the contradiction, completing every terminal branch.
\end{proof}

\paragraph{Scope and remaining boundary.}
This removes tensor indecomposability from D-TD36 at the explicit
larger threshold $n\ge2^{30}$.  The whole commuting-field model
may be a nontrivial tensor product, and the actual quotient need
not lift to any normal alternating subgroup of the original $H$.
The only normal perfect subgroup constructed above lives in the
selected finite coordinate lift where it is used.
Absolute primitivity of the original commuting-field model is
still required.  Absolutely imprimitive defining modules and
kernel-field sign interfaces not satisfying (HPT.1)--(HPT.2)
are not eliminated.  No original characteristic-zero field,
polarization or successful source is constructed.
